% Source: https://artofproblemsolving.com/community/c3h2275279_500th_post
% Stable identity: aops-topic-2275279-q01
The semiperimeter is $s=\dfrac{6+9+12}{2}=\dfrac{27}{2}$, and by Heron's formula the area is
\[K=\sqrt{\dfrac{27}{2}\cdot\dfrac{15}{2}\cdot\dfrac{9}{2}\cdot\dfrac{3}{2}}=\dfrac{\sqrt{10935}}{4}=\dfrac{27\sqrt{15}}{4}.\]
The incentre $I$ splits the triangle into three triangles with bases $6$, $9$, $12$ and common height $r$, so $K=rs$ and
\[r=\dfrac Ks=\dfrac{27\sqrt{15}}{4}\cdot\dfrac{2}{27}=\boxed{\dfrac{\sqrt{15}}{2}}\approx1.94.\]

\begin{asy}
usepackage("amsmath");
size(5.5cm);
defaultpen(fontsize(10pt));
pair B=(0,0), C=(12,0); real ax=99/24; pair A=(ax,sqrt(36-ax^2));
pair I=(12*A+9*B+6*C)/27; real r=sqrt(15)/2;
draw(A--B--C--cycle,linewidth(0.9)); draw(circle(I,r),blue); draw(I--(I.x,0),blue);
dot(I); label("$I$",I,N,blue); label("$r$",(I.x,r/2),E,blue);
label("$A$",A,N); label("$B$",B,SW); label("$C$",C,SE);
label("$6$",(A+B)/2,NW); label("$9$",(A+C)/2,NE); label("$12$",(9,0),S);
\end{asy}
